\documentclass[11pt]{article}
\usepackage[margin=1in]{geometry}
\usepackage{amsmath,amssymb}
\usepackage{graphicx}
\usepackage{booktabs}
\usepackage{hyperref}
\usepackage{natbib}
\usepackage{setspace}

\title{Biphasic V1 Population Dynamics Underlying Perceptual Similarity Masking Revealed by Voltage-Sensitive Dye Imaging in Behaving Macaques}
\author{Author A.$^{1}$, Author B.$^{1,2}$, Author C.$^{2}$ \\
\small $^{1}$Department of Neuroscience, University of Example \\
\small $^{2}$Center for Visual Neuroscience, Example Institute}
\date{}

\begin{document}
\maketitle

\begin{abstract}
Similarity masking---the reduced detectability of a target as it becomes more similar to its background---is a fundamental property of visual perception that underlies natural camouflage. Despite extensive study of contextual modulation in primary visual cortex (V1), the neural population dynamics that give rise to orientation-dependent masking during active behavior remain unknown. Here we used voltage-sensitive dye imaging (VSDI) to record V1 population activity in two macaque monkeys performing a target detection task with systematically varied target-background orientation similarity across five contrast combinations. Behavioral data confirmed that detection sensitivity (d$'$) was lowest when the target and background orientations matched, replicating the hallmark of similarity masking. Neural analysis revealed a biphasic temporal profile at both retinotopic and columnar spatial scales: an initial transient phase in which similar backgrounds paradoxically enhanced target responses, followed by a sustained phase in which similar backgrounds suppressed responses in a pattern consistent with behavioral masking. Critically, columnar-scale signals---reflecting orientation-selective competition among V1 columns---correlated with behavioral masking earlier and more robustly than retinotopic-scale signals, with integrated columnar responses positively correlated with behavioral d$'$ across all five contrast conditions. Decomposition into 12 orientation-selective clusters revealed that sustained suppression was driven by reduced activation of target-aligned neurons while orthogonal neurons were relatively unaffected. A delayed divisive normalization model qualitatively reproduced the biphasic dynamics, suggesting that the temporal lag of a normalization pool explains both the initial paradoxical enhancement and the subsequent suppression. These results establish that V1 population dynamics at the columnar scale carry the critical neural correlate of similarity masking during natural visual behavior.
\end{abstract}

\section{Introduction}

The ability to detect a target amid a structured background is a ubiquitous perceptual challenge with direct ecological relevance: camouflage exploits similarity masking to render prey invisible to predators. In human psychophysics, similarity masking manifests as a monotonic decrease in detection performance as the orientation difference between a target grating and a surrounding mask approaches zero \citep{foley1994,sebastian2017}. Despite decades of research on contextual modulation in V1 \citep{hubel1962,blakemore1970,cavanaugh2002,angelucci2017}, the spatiotemporal dynamics of V1 population responses that underlie this fundamental perceptual phenomenon during active visual behavior have not been directly characterized.

Prior work on masking has focused primarily on contrast gain control mechanisms \citep{heeger1992,carandini2012}, demonstrating that surround stimuli suppress V1 neuronal responses through divisive normalization. However, these studies typically employed anesthetized preparations or passive fixation paradigms and did not dissect the temporal evolution of masking signals at different spatial scales within the V1 population. Furthermore, the critical question of whether orientation-specific interactions at the scale of orientation columns versus the broader retinotopic scale drive behavioral masking has not been addressed.

Here we combined voltage-sensitive dye imaging (VSDI) through a cranial chamber with transparent artificial dura \citep{grinvald1999} with a target detection task to probe the neural basis of similarity masking in two behaving macaques. VSDI provides millisecond temporal resolution over millimeters of cortex, enabling simultaneous analysis at retinotopic (target location) and columnar (orientation preference) scales. We report a previously unknown biphasic temporal profile of V1 population responses and demonstrate that columnar-scale signals provide the dominant neural correlate of behavioral masking.

\section{Methods}

\subsection{Subjects and Recording Setup}

Two male \textit{Macaca mulatta} were implanted with cranial chambers over V1, fitted with transparent artificial dura for chronic optical access. VSDI was performed through an 8$\times$8~mm imaging window. Retinotopic mapping confirmed that the imaged region corresponded to approximately 3~degrees eccentricity and that V2 was completely hidden within the lunate sulcus.

\subsection{Behavioral Task}

Monkeys performed a target detection task. On each trial, a 4-degree background grating (4~cpd, cosine-centered, raised-cosine-masked) was presented at approximately 3~degrees eccentricity. On 50\% of trials, a small horizontal Gabor target (4~cpd, 0.33-degree FWHM envelope) was added at the center of the background. Background orientation varied randomly from trial to trial between 0~degrees (matching the horizontal target) and 90~degrees (orthogonal). Monkeys indicated target presence with a saccade to the target location and target absence by maintaining fixation.

\begin{table}[ht]
\centering
\caption{Contrast conditions tested in the detection task. Target contrast (T) and background contrast (Bg) are expressed as percentages.}
\label{tab:contrast}
\begin{tabular}{lcc}
\toprule
Condition & Target contrast (\%) & Background contrast (\%) \\
\midrule
T12 Bg7  & 12 & 7 \\
T24 Bg12 & 24 & 12 \\
T12 Bg12 & 12 & 12 \\
T24 Bg24 & 24 & 24 \\
T06 Bg7  & 6  & 7 \\
\bottomrule
\end{tabular}
\end{table}

Five combinations of target and background contrast were tested (Table~\ref{tab:contrast}). Behavioral performance was quantified using signal detection theory \citep{green1966}: d$' = z(\text{hit rate}) - z(\text{false alarm rate})$.

\subsection{Neural Analysis}

Neural analysis was performed at two spatial scales. At the \textbf{retinotopic scale}, a Gaussian spatial template matched to the target's retinotopic profile was applied to extract the amplitude of the target-evoked VSD response. The template was doubly whitened for noise and target-absent baseline.

At the \textbf{columnar scale}, an orientation-selective template was computed from each experiment's V1 orientation map. This template extracted the relative strength of activity in columns tuned to the target orientation (0~degrees) versus columns tuned to orthogonal orientations. Positive values indicated stronger activation of target-tuned columns.

For finer analysis, each pixel in the imaging field was assigned to one of 12 equally spaced orientation clusters based on its preferred orientation from the orientation map. Population tuning curves were constructed by computing the mean VSD response within each cluster, windowed to the retinotopic profile of the target.

\subsection{Neural-Behavioral Correlation}

At each time point during the trial, we computed the correlation between the neural template response (as a function of background orientation) and the behavioral d$'$ (as a function of background orientation). This time-resolved correlation tracked when V1 population signals began to predict behavioral masking.

\subsection{Normalization Model}

A delayed divisive normalization model was used to test whether a simple gain control computation could account for the observed biphasic dynamics \citep{carandini2012,reynoldsheeger2009}. The model comprised an excitatory drive that responded immediately to the stimulus and a normalization pool that built up with a temporal delay.

\section{Results}

\subsection{Behavioral Similarity Masking}

Both monkeys exhibited clear similarity masking across all five contrast conditions. Detection d$'$ was lowest when the background orientation matched the target (0~degrees) and highest when the background was orthogonal (90~degrees) or absent (uniform gray). Performance on uniform backgrounds exceeded that on any grating background, indicating a general masking effect of textured surrounds. The orientation-dependent component of masking was well described by an inverted Gaussian function at each contrast level.

\begin{table}[ht]
\centering
\caption{Summary of behavioral similarity masking across contrast conditions. All conditions showed the characteristic pattern of lowest d$'$ at matched orientation.}
\label{tab:behavior}
\begin{tabular}{lccc}
\toprule
Condition & d$'$ at 0$^{\circ}$ & d$'$ at 90$^{\circ}$ & d$'$ (uniform) \\
\midrule
T12 Bg7  & Lowest & Higher & Highest \\
T24 Bg12 & Lowest & Higher & Highest \\
T12 Bg12 & Lowest & Higher & Highest \\
T24 Bg24 & Lowest & Higher & Highest \\
T06 Bg7  & Lowest & Higher & Highest \\
\bottomrule
\end{tabular}
\end{table}

Reaction times on hit trials showed a complex, sometimes non-monotonic, dependence on target-background orientation similarity.

\subsection{Biphasic Dynamics at the Retinotopic Scale}

VSD responses at the retinotopic scale revealed a striking biphasic temporal pattern. During the initial transient phase (approximately 50--100~ms post-stimulus), target responses were paradoxically enhanced by similar backgrounds: the retinotopic template response was larger when the background matched the target orientation compared to orthogonal backgrounds. During the subsequent sustained phase (approximately 100--250~ms post-stimulus), this pattern reversed: similar backgrounds now suppressed the target response relative to orthogonal backgrounds, consistent with the direction of behavioral masking.

The time-resolved neural-behavioral correlation reflected this reversal. Early after stimulus onset, the correlation between neural sensitivity and behavioral d$'$ across orientations was negative (anti-correlated), indicating that the initial enhancement by similar backgrounds opposed the behavioral masking pattern. The correlation then reversed sign, becoming positive during the sustained phase. However, when temporally integrated over the full response epoch, retinotopic signals were often uncorrelated or anti-correlated with behavioral performance, suggesting that this spatial scale alone does not fully account for masking.

\subsection{Columnar-Scale Signals Dominate Behavioral Correlates}

At the columnar scale, the biphasic pattern was more pronounced. The sign reversal from anti-correlation to positive correlation between neural and behavioral masking occurred earlier than at the retinotopic scale. Most importantly, temporally integrated columnar responses were positively correlated with behavioral d$'$ across all five contrast conditions---a consistency that was absent at the retinotopic scale.

\begin{table}[ht]
\centering
\caption{Comparison of retinotopic and columnar scale signals in predicting behavioral masking.}
\label{tab:scales}
\begin{tabular}{lcc}
\toprule
Feature & Retinotopic scale & Columnar scale \\
\midrule
Biphasic dynamics & Present & Present (more pronounced) \\
Positive correlation onset & Later & Earlier \\
Consistency across contrasts & Variable & All conditions positive \\
Integrated response--behavior & Often uncorrelated & Positively correlated \\
\bottomrule
\end{tabular}
\end{table}

This dissociation (Table~\ref{tab:scales}) indicates that behavioral masking is driven primarily by the balance of activity across orientation-selective V1 columns, rather than by the overall amplitude of the target's retinotopic response.

\subsection{Orientation-Selective Decomposition}

Decomposing the VSD response into 12 orientation-selective clusters revealed the source of the sustained columnar signal. During the initial transient, population tuning curves peaked at the background orientation, reflecting the dominant feedforward drive. During the sustained phase, however, target-aligned neurons (those whose preferred orientation matched the target) showed reduced activation when the background matched the target orientation, while neurons tuned to orthogonal orientations were relatively unaffected. This selective suppression of target-aligned columns by similar backgrounds drives the columnar masking signal and explains why target detection becomes more difficult when the background is similar.

\subsection{Correct-Trial Analysis}

Re-analysis restricted to correct trials (hits and correct rejections) confirmed that the biphasic dynamics were not an artifact of decision-related or attention-related top-down modulation. The temporal profile and sign reversal were preserved in the correct-trial subset, supporting a primarily feedforward mechanism.

\subsection{Delayed Divisive Normalization Accounts for Biphasic Dynamics}

The delayed divisive normalization model qualitatively reproduced the biphasic pattern. In the model, the early response was dominated by excitatory feedforward drive, which is stronger when the background and target activate overlapping neural populations (similar orientations). Before the normalization pool fully engages, similar backgrounds therefore enhance target responses. Once normalization builds up with its characteristic delay, the suppressive interaction dominates: similar backgrounds produce stronger normalization of target-aligned neurons, reducing the net target response. The sign reversal timing in the model depended on the normalization delay parameter and qualitatively matched the observed reversal in the data.

\section{Discussion}

We have identified a biphasic temporal profile of V1 population responses that underlies perceptual similarity masking in behaving macaques. The initial paradoxical enhancement of target responses by similar backgrounds, followed by sustained suppression, can be explained by a delayed divisive normalization mechanism \citep{carandini2012}. Critically, the columnar scale of V1 orientation maps provides the dominant neural correlate of behavioral masking, with integrated columnar signals predicting detection performance across all tested contrast conditions.

\subsection{Columnar Versus Retinotopic Contributions}

Our finding that columnar-scale signals provide a more robust and consistent predictor of behavioral masking than retinotopic-scale signals was unexpected. Standard models of target detection emphasize the retinotopic target response---the signal at the target's spatial location---as the primary determinant of detectability. Our results suggest that the brain instead reads out a comparison between target-tuned and orthogonal columns within the target's retinotopic neighborhood. This columnar readout is inherently more sensitive to orientation-specific masking because it directly reflects the competition between orientation representations.

\subsection{Relationship to Normalization Models}

The delayed divisive normalization framework \citep{reynoldsheeger2009,schwartz2001} provides an elegant account of the biphasic dynamics. The key parameter is the temporal delay of the normalization pool. During the lag period, the excitatory drive from both target and background is un-normalized, and similar backgrounds boost the population response at the target location. Once normalization engages, the suppressive effect of background similarity dominates. This temporal structure predicts that very brief targets would escape similarity masking---a prediction consistent with psychophysical observations of reduced masking at short stimulus durations.

\subsection{Implications for Camouflage and Natural Vision}

Similarity masking is the primary mechanism underlying visual camouflage in nature. Our results suggest that camouflaged animals exploit not just the initial feedforward response (where similarity paradoxically enhances visibility) but the sustained normalization-driven suppression that follows. The time course we observed---with suppression emerging around 100~ms---is consistent with the duration of typical fixations during natural viewing, suggesting that similarity masking operates on a behaviorally relevant timescale.

\subsection{Limitations}

Our study imaged only V1 and cannot speak to the contributions of downstream areas such as V2 and V4 to similarity masking. The normalization model is descriptive rather than mechanistic; the biological implementation of the delayed normalization pool likely involves both local inhibitory circuits and feedback from higher areas \citep{angelucci2017}. Additionally, our target was always horizontal, so we cannot distinguish target-specific effects from orientation-specific biases in V1 columnar architecture.

\section{Conclusion}

We have demonstrated that V1 population responses exhibit a biphasic temporal profile during similarity masking in behaving primates. An initial paradoxical enhancement by similar backgrounds is followed by sustained suppression, with columnar-scale orientation signals providing the critical neural correlate of behavioral detection performance. A delayed divisive normalization model accounts for this temporal structure. These findings advance our understanding of how V1 population dynamics support visual perception during natural behavior.

\bibliographystyle{plainnat}
\bibliography{references}

\end{document}
